% This template has been taken and adapted from ICML 2021 submission file available under Creative Commons CC BY 4.0.

% The main.tex ICML 2021 submission file contains the following attribution:

% "This document was modified from the file originally made available by
% Pat Langley and Andrea Danyluk for ICML-2K. This version was created
% by Iain Murray in 2018, and modified by Alexandre Bouchard in
% 2019 and 2021. Previous contributors include Dan Roy, Lise Getoor and Tobias
% Scheffer, which was slightly modified from the 2010 version by
% Thorsten Joachims & Johannes Fuernkranz, slightly modified from the
% 2009 version by Kiri Wagstaff and Sam Roweis's 2008 version, which is
% slightly modified from Prasad Tadepalli's 2007 version which is a
% lightly changed version of the previous year's version by Andrew
% Moore, which was in turn edited from those of Kristian Kersting and
% Codrina Lauth. Alex Smola contributed to the algorithmic style files."

% Link to the original template: https://icml.cc/Conferences/2021/StyleAuthorInstructions

% The adaptation of the ICML 2021 submission file was done at the EPFL and Idiap Research Institute by Alina Elena Baia, Darya Baranouskaya and Olena Hrynenko (equal contribution).

\documentclass{article}

% importing the style files
\input{source/preamble}
\usepackage[accepted]{source/icml2021}
\usepackage{caption}
\usepackage{titlesec}
\usepackage{siunitx}

\titlespacing*{\section}{0pt}{6pt}{3pt}
\titlespacing*{\subsection}{0pt}{4pt}{2pt}


%%%%%%%%% ADD YOUR GROUP NUMBER %%%%%%%%%
\icmltitlerunning{EE-559 Deep Learning: Group Mini-Project, Group 44}

\begin{document}

\twocolumn[
%%%%%%%%% ADD TITLE OF YOUR GROUP MINI-PROJECT %%%%%%%%%
\vspace{-35px}
\icmltitle{Detect and Explain: RAC and LLM explainability on Implicit hate speech.}

%%%%%%%%% ADD NAMES OF THE GROUP MEMBERS %%%%%%%%%
\begin{icmlauthorlist}
\icmlauthor{Pierre\ Bernadet}{ch} \quad
\icmlauthor{Alexandre\ Potocnik}{ch} \quad
\icmlauthor{Julie\ Michallat}{ch} \\
\end{icmlauthorlist}

% adding affiliation information
\icmlaffiliation{ch}{}
\centering
%%%%%%%%% ADD YOUR GROUP NUMBER %%%%%%%%%
Group 44
\vskip 0.1in
]

% command for printing affiliation
%%\printAffiliationsAndNotice{}

%%%%%%%%% ADD ABSTRACT %%%%%%%%%
\begin{abstract}
Implicit hate speech — relying on coded language, dog whistles, and irony — remains significantly harder to detect than its explicit counterpart, while content moderation regulations require justifying removal decisions. We propose a Retrieval-Augmented Classifier (RAC) that grounds fine-tuned transformer encoders in two complementary FAISS indexes: labeled precedents from three benchmark datasets and expert knowledge bases. The Example and Full indexes consistently improve over no-index baselines, with best gains of $+$1.4\,pp F1 on IHC (F1\,=\,0.811), $+$2.7\,pp F1 on ISHate (F1\,=\,0.906), and $+$2.6\,pp F1 on Vicomtech (F1\,=\,0.854). On-demand LLaMA~3.3 explanations anchored to retrieved evidence address explainability requirements without involving an LLM in the classification decision.

%%%%%%%%% ADD KEYWORDS %%%%%%%%%
% Make sure to keep keywords on a new line.
\textbf{Keywords:} implicit hate speech, retrieval-augmented classification, LLM explainability

\end{abstract}

%%%%%%%%% ADD INTRODUCTION SECTION %%%%%%%%%
\vspace{-20pt}
\section{Introduction}
\label{sec:intro}

Hate speech detection has made substantial progress with fine-tuned transformer encoders, achieving F1 scores above 0.90 on explicit content \cite{saleh2021bert}. Implicit hate speech, however, relies on coded language, dog whistles, and irony to convey hatred without surface-level slurs, making it significantly harder to detect \cite{elsherief2021latent, ocampo2023indepth}. Unlike explicit forms, it cannot be identified from lexical patterns alone; understanding it requires broader cultural and contextual knowledge that standard classifiers lack. Moreover, the EU Digital Services Act now requires platforms to justify content removal decisions, yet existing encoder-based classifiers offer no explanation mechanisms.

We address both challenges with a RAC: retrieval grounds detection in labeled precedents and expert knowledge, while providing a natural anchor for human-readable explanation. Concretely, we contribute (1) a RAC pipeline combining a Sentence-BERT retriever \cite{reimers2019sbert} with a dual FAISS index (labeled examples + expert glossaries) feeding a fine-tuned encoder; (2) a systematic evaluation across 3 encoders, 3 datasets, and 3 index configurations; (3) a hyperparameter study ($k$, $\tau$) demonstrating low retrieval sensitivity; and (4) on-demand LLaMA~3.3 explanations grounded in retrieved evidence, decoupled from the classification decision.


%%%%%%%%% ADD RELATED WORK SECTION %%%%%%%%%
\section{Related Work}
\label{sec:related_work}

\textbf{Implicit hate speech benchmarks and encoder-based classifiers.}
ElSherief et al.~\cite{elsherief2021latent} introduce IHC and established the first benchmarks on this task.
De~Gibert et al.~\cite{degibert2018hate} note that the Vicomtech corpus's single-source origin limits domain generalisability.
Ocampo et al.~\cite{ocampo2023indepth} show that state-of-the-art encoders ``fail to detect implicit and subtle content'' across seven benchmarks.

\textbf{RAG for text classification and content moderation.}
Policy-grounded RAG systems applied to hate speech~\cite{mnassri2025rag} exist but rely on a generative LLM as the final decision-maker, increasing inference and limiting training scalability.

\textbf{Current state of the art on IHC.}
Cheremetiev et al.~\cite{cheremetiev2025specializing} achieve F1-macro of 81.68 on IHC with large embedding models requiring up to 62.4\,GB of GPU memory and no explanation mechanism.


%%%%%%%%% ADD METHOD SECTION %%%%%%%%%
\section{Method}
\label{sec:method}

\subsection{Datasets}
We use three implicit hate speech benchmarks unified into a binary task: \texttt{not\_hate} (0) vs.\ \texttt{hate} (1). \textbf{IHC}~\cite{elsherief2021latent} contains tweets with explicit and implicit hate labels, collapsed into a single \texttt{hate} class with a 90/10 train/test split. \textbf{ISHate}~\cite{ocampo2023indepth} provides examples with pre-made splits; all hate sub-types are collapsed to binary. \textbf{Vicomtech}~\cite{degibert2018hate} consists of posts with pre-made binary splits from a white-supremacist forum called Stormfront.

\subsection{Expert Knowledge Bases}
We aggregate three hate speech resources: the \textbf{ADL}~\cite{adl} glossary of hate symbols and extremist slang, the \textbf{SPLC}~\cite{splc} annotated extremist terminology, and \textbf{HateBase}~\cite{hatebase}, a crowdsourced multilingual slur repository. After scraping, chunking, and deduplication, these yield 40,951 entries forming the Knowledge index.

\subsection{Vector Index Construction}
All chunks are encoded with \textbf{Sentence-BERT}~\cite{reimers2019sbert} and stored in L2-normalised \textbf{FAISS} indexes. We construct three indexes: \textbf{Example (E)}, label-prefixed chunks from training splits of all three datasets; \textbf{Knowledge (K)}, chunks from the expert glossaries; and \textbf{Full (E+K)}, their union. During training, the input is masked to prevent trivial self-retrieval.

\subsection{Augmented Input Format}
The retriever fetches the $k$ nearest chunks exceeding cosine similarity threshold $\tau$. The augmented input takes the form $\texttt{[query]} \;\texttt{[SEP]}\; \texttt{[}l_1\texttt{]}\; n_1 \;\texttt{[SEP]}\; \texttt{[}l_2\texttt{]}\; n_2 \;\texttt{[SEP]}\; \cdots$, where $n_i$ is the $i$-th retrieved neighbour and $l_i \in \{\texttt{hate},\, \texttt{not hate}\}$ its label prefix. Retrieved neighbours act as labelled few-shot anchors for the classification decision.

\subsection{Classifier Fine-tuning}
We fine-tune three encoders: \textbf{BERT},  \textbf{HateBERT} (110M parameters) and \textbf{RoBERTa} (125M parameters), pre-trained on abusive language corpora. Each is fine-tuned with a binary classification head, for a standard recommendation of 3 epochs (lr\,=\,$2{\times}10^{-5}$, batch size 16, seed 42). To ensure full reproducibility, the random seed is fixed before each classifier head initialisation, independently of the training loop seed.

\subsection{Hyperparameter Tuning}
We grid-search $k \in \{3, 5, 10\}$ and $\tau \in \{0.3, 0.4, 0.5, 0.6\}$ on IHC with the Full index (Tab.~\ref{tab:hp_tuning}, Appendix). Optimal values are applied to all subsequent experiments, as hyperparameter sensitivity has been found to be low.

\subsection{On-demand Explanation}
At inference time, the augmented input and the encoder's prediction are passed to \textbf{LLaMA~3.3} to generate a grounded, human-readable justification. The LLM is not involved in the classification decision and is not fine-tuned; it synthesizes retrieved evidence into a natural-language explanation on demand.

\vspace{-2px}
\begin{figure}[h]
    \centering
    \includegraphics[width=\columnwidth]{model_overview3.png}
    \captionsetup{font=footnotesize}
    \caption{Overview of the RAC and LLM pipeline.}
    \label{fig:pipeline}
\end{figure}

%%%%%%%%% ADD VALIDATION SECTION %%%%%%%%%
\section{Validation}
\label{sec:validation}

\subsection{Experimental Setup}
We evaluate 27 RAC models (3 encoders $\times$ 3 indexes $\times$ 3 datasets) plus 9 no-index baselines, using macro F1, precision, and recall. Hyperparameters are fixed to values from Sec.~3.6. A single seed (42) is used throughout. To estimate variance, RoBERTa(E+K) was additionally evaluated across 4 seeds (Tab.~\ref{tab:multiseed}, Appendix): IHC ($0.800\pm0.003$) and ISHate ($0.900\pm0.005$) are stable, while Vicomtech shows higher variance ($0.831\pm0.016$), consistent with its small training set. Tab.~\ref{tab:hp_tuning} (Appendix) confirms F1 variation across the full HP grid remains below 2.5 pp per model, validating low retrieval sensitivity to $k$ and $\tau$.



\subsection{Main Results}
Tab.~\ref{tab:results} reports macro F1 for all configurations; full precision and recall are in Tab.~\ref{tab:full_results} (Appendix).

\textbf{Retrieval consistently improves over baselines for the Example and Full indexes.} The Knowledge index, however, underperforms the no-index baseline in nearly all settings, suggesting that the automatically-scraped expert glossaries introduce noise rather than useful signal (see also Sec.~\ref{sec:limits}).

\textbf{IHC and ISHate.} On IHC, RoBERTa(E) achieves the best result of F1\,=\,0.811 ($+$1.4 pp over its no-index baseline). On ISHate, RoBERTa(E+K) achieves F1\,=\,0.906 ($+$2.7 pp), with the Full index outperforming the Example-only index. Retrieved labelled examples act as effective few-shot anchors grounding the encoder in semantically similar precedents.

\textbf{Vicomtech.} On Vicomtech (1,914 training examples), RoBERTa(E+K) achieves F1\,=\,0.854 ($+$2.6 pp over baseline), with RoBERTa(E)\,=\,0.845 close behind. The Full index outperforms the Example-only index, confirming that combining both sources provides consistent additional signal even on small datasets.

\textbf{RoBERTa consistently outperforms BERT} across all datasets and index types.

\textbf{Our models substantially outperform published baselines.} Tab.~\ref{tab:ihc_comparison} and Tab.~\ref{tab:ishate_taskA} (Appendix) replicate the exact setups of ElSherief et al.~\cite{elsherief2021latent} and Ocampo et al.~\cite{ocampo2023indepth}: RoBERTa (RAC) reaches F1\,=\,0.795 vs.\ BERT+Aug at 0.704 on IHC, and baselines are competitive on ISHate Task A except on Implicit HS, where the binary retrieval signal conflates all three classes.

\begin{table}[h]
\centering
{\footnotesize
\begin{tabular}{lccc}
\toprule
\textbf{Model (Index)} & \textbf{IHC} & \textbf{ISHate} & \textbf{Vicomtech} \\
\midrule
BERT (No-index)     & 0.786 & 0.872 & 0.814 \\
BERT (E)            & 0.782 & 0.892 & 0.842 \\
BERT (K)            & 0.768 & 0.872 & 0.805 \\
BERT (E+K)          & 0.786 & 0.891 & 0.825 \\
\midrule
HateBERT (No-index) & 0.787 & 0.881 & 0.822 \\
HateBERT (E)        & 0.779 & 0.900 & 0.834 \\
HateBERT (K)        & 0.769 & 0.873 & 0.816 \\
HateBERT (E+K)      & 0.778 & 0.898 & 0.840 \\
\midrule
RoBERTa (No-index)  & 0.797 & 0.879 & 0.828 \\
RoBERTa (E)         & \textbf{0.811} & 0.903 & 0.845 \\
RoBERTa (K)         & 0.792 & 0.888 & 0.808 \\
RoBERTa (E+K)       & 0.800 & \textbf{0.906} & \textbf{0.854} \\
\bottomrule
\end{tabular}}
\captionsetup{font=footnotesize}
\caption{Macro F1 results. E\,=\,Example index, K\,=\,Knowledge index, E+K\,=\,Full index. Bold indicates best per column. Full results (Prec./Rec.) in Tab.~\ref{tab:full_results} (Appendix).}
\label{tab:results}
\vspace{-10pt}
\end{table}


\subsection{Ablation: Removing IHC from the Retrieval Index}
We train and evaluate on IHC while excluding all IHC training examples from the retrieval index (Tab.~\ref{tab:noihc_ablation}, Appendix). RoBERTa achieves F1\,=\,0.801 with the no-IHC index, compared to 0.811 with the full Example index. This difference falls within the seed variance of RoBERTa(E+K) on IHC ($0.800\pm0.003$, Tab.~\ref{tab:multiseed}), confirming that the retriever generalises via semantic similarity rather than exploiting in-distribution shortcuts.


\subsection{Qualitative Analysis}
Fig.~\ref{fig:rac-example} illustrates the RAC pipeline on an input containing the antisemitic acronym \textit{ZOG} (Zionist Occupation Government). The retriever surfaces a HateBase definition of ZOG as the top match alongside two labelled hate examples. HateBERT(E+K) assigns \texttt{hate}; LLaMA~3.3 synthesises the retrieved evidence into a grounded explanation flagging the antisemitic framing.


\begin{figure}[h]
\centering
\setlength{\fboxsep}{6pt}
\resizebox{0.95\columnwidth}{!}{%
\fbox{%
\begin{minipage}{\dimexpr\columnwidth-2\fboxsep-2\fboxrule\relax}\small
  \textbf{Input:} \textit{``You will be drafted to defend zog''}\\[1pt]
  {\footnotesize Label: hate \;|\; Prediction: hate \;|\; Confidence: 99.2\%}

  \vspace{-5pt}\rule{\linewidth}{0.3pt}\vspace{0pt}

  \textsc{Retrieved Passages}\\[1pt]
  \begin{tabular}{@{} r l p{0.78\linewidth} @{}}
    \# & Score & Passage \\\hline\\[-7pt]
    1 & 0.445 & [Hate] zog: Zionist Occupation Government, an organization hypothesized by white supremacists to rule the world. Target category: Judaism. \\[1pt]
    2 & 0.427 & [Hate] this is antifa . zog will take anyone to fight for greater israel \\[1pt]
    3 & 0.419 & [Hate] Maximum resistance towards Zog! \\[0pt]
  \end{tabular}

  \rule{\linewidth}{0.3pt}\vspace{2pt}

  \textsc{Model Explanation}\\[1pt]
  The content was flagged as hate speech due to the reference to \textit{zog}, a term associated with white supremacist ideology and anti-Semitic conspiracy theories. This term is often used to express hatred towards the Jewish community.
\end{minipage}%
}%
}
\captionsetup{font=footnotesize}
\caption{Example of RAC-augmented hate speech detection. Retrieved passages provide
contextual evidence guiding the model explanation.}
\label{fig:rac-example}
\vspace{-10pt}
\end{figure}





\subsection{Computational Cost}
Our RoBERTa(E+K) requires 0.88\,GB of GPU memory and achieves a slightly lower F1 of 0.811 compared to NV-Embed (62.4\,GB ; F1 0.816). 
This 70$\times$ memory reduction makes the framework substantially more practical for large-scale deployment under DSA cost constraints.



\subsection{Limitations}
\label{sec:limits}
\textbf{Annotation subjectivity.} Implicit hate speech lacks a universal definition, leading to annotator disagreement.

\textbf{Knowledge base scraping quality.}
The expert glossaries were scraped automatically, introducing two quality issues. First, long entries were split into overlapping chunks using a fixed sliding-window, breaking definitions mid-sentence. Second, HateBase entries were often scraped without human-written definitions, retaining only structured metadata, yielding semantically vacuous retrieved fragments. This likely explains why the Knowledge index consistently underperforms the Example index.

\textbf{Explanation faithfulness.} As classifier and explainer are decoupled, LLaMA~3.3 explanations may reference evidence that did not drive the classification score.


%%%%%%%%% ADD CONCLUSION SECTION %%%%%%%%%
\section{Conclusion}
\label{sec:conclusion}

We applied Retrieval-Augmented Classification to implicit hate speech detection, grounding fine-tuned encoders in labeled precedents and expert knowledge bases. The Example and Full indexes consistently improve over no-index baselines across all three datasets, with gains up to $+$1.4 pp on IHC, $+$2.7 pp on ISHate, and $+$2.6 pp on Vicomtech. On-demand LLaMA~3.3 explanations address regulatory explainability requirements without involving a generative model in the classification decision. Future work includes multilingual extension, entry-aware chunking to resolve the HateBase scraping quality issue, and joint retriever--classifier fine-tuning for maximally discriminative retrieval.


%%%%%%%%% REFERENCES SECTION %%%%%%%%%
{\tiny
\setlength{\bibsep}{2pt}
\bibliographystyle{ieeetr}
% This is the link to the refereces file.
\bibliography{main}
}



\appendix
\onecolumn

\section{Full Results}

\begin{center}
\captionof{table}{RAC fine-tuning results (macro F1, Precision, Recall). E\,=\,Example index, K\,=\,Knowledge index, E+K\,=\,Full index. Bold indicates best per column.}
\label{tab:full_results}
\resizebox{\textwidth}{!}{%
\begin{tabular}{lccccccccc}
\toprule
 & \multicolumn{3}{c}{\textbf{IHC}} & \multicolumn{3}{c}{\textbf{ISHate}} & \multicolumn{3}{c}{\textbf{Vicomtech}} \\
\cmidrule(lr){2-4} \cmidrule(lr){5-7} \cmidrule(lr){8-10}
\textbf{Model (Index)} & F1 & Prec. & Rec. & F1 & Prec. & Rec. & F1 & Prec. & Rec. \\
\midrule
BERT (No-index)     & 0.786 & 0.791 & 0.783 & 0.872 & 0.873 & 0.871 & 0.814 & 0.814 & 0.814 \\
BERT (E)            & 0.782 & 0.789 & 0.777 & 0.892 & 0.890 & 0.894 & 0.842 & 0.852 & 0.843 \\
BERT (K)            & 0.768 & 0.773 & 0.764 & 0.872 & 0.872 & 0.873 & 0.805 & 0.806 & 0.805 \\
BERT (E+K)          & 0.786 & 0.792 & 0.781 & 0.891 & 0.890 & 0.893 & 0.825 & 0.835 & 0.826 \\
\midrule
HateBERT (No-index) & 0.787 & 0.793 & 0.783 & 0.881 & 0.881 & 0.882 & 0.822 & 0.823 & 0.822 \\
HateBERT (E)        & 0.779 & 0.784 & 0.775 & 0.900 & 0.898 & 0.902 & 0.834 & 0.843 & 0.835 \\
HateBERT (K)        & 0.769 & 0.773 & 0.766 & 0.873 & 0.874 & 0.873 & 0.816 & 0.817 & 0.816 \\
HateBERT (E+K)      & 0.778 & 0.783 & 0.774 & 0.898 & 0.896 & 0.900 & 0.840 & 0.851 & 0.841 \\
\midrule
RoBERTa (No-index)  & 0.797 & 0.802 & 0.793 & 0.879 & 0.877 & 0.882 & 0.828 & 0.829 & 0.828 \\
RoBERTa (E)         & \textbf{0.811} & \textbf{0.820} & \textbf{0.805} & 0.903 & 0.900 & 0.907 & 0.845 & 0.847 & 0.845 \\
RoBERTa (K)         & 0.792 & 0.799 & 0.788 & 0.888 & 0.885 & 0.891 & 0.808 & 0.808 & 0.808 \\
RoBERTa (E+K)       & 0.800 & 0.806 & 0.795 & \textbf{0.906} & \textbf{0.904} & \textbf{0.909} & \textbf{0.854} & \textbf{0.854} & \textbf{0.854} \\
\bottomrule
\end{tabular}}
\end{center}

\section{Hyperparameter Tuning \& Seed Variance}

\begin{center}
\captionof{table}{Macro F1 over the HP grid ($k$, $\tau$) on IHC with the Full index. Best per model in bold.}
\label{tab:hp_tuning}
\resizebox{0.65\textwidth}{!}{%
\begin{tabular}{llcccc}
\toprule
\textbf{Model} & $k$ & $\tau{=}0.3$ & $\tau{=}0.4$ & $\tau{=}0.5$ & $\tau{=}0.6$ \\
\midrule
\multirow{3}{*}{BERT}     & 3  & 0.777 & 0.779 & 0.787 & 0.785 \\
                           & 5  & 0.775 & 0.775 & 0.778 & 0.779 \\
                           & 10 & 0.780 & 0.776 & 0.783 & \textbf{0.788} \\
\midrule
\multirow{3}{*}{HateBERT} & 3  & 0.777 & \textbf{0.781} & 0.777 & 0.775 \\
                           & 5  & 0.771 & 0.768 & 0.775 & 0.773 \\
                           & 10 & 0.759 & 0.768 & 0.775 & 0.775 \\
\midrule
\multirow{3}{*}{RoBERTa}  & 3  & 0.800 & 0.806 & 0.802 & 0.786 \\
                           & 5  & \textbf{0.809} & 0.799 & 0.802 & 0.797 \\
                           & 10 & 0.801 & 0.801 & 0.801 & 0.794 \\
\bottomrule
\end{tabular}}
\end{center}

\vspace{0.5em}

\begin{center}
\captionof{table}{RoBERTa\,(E+K) macro F1 across 4 seeds. High variance on Vicomtech reflects its small training set.}
\label{tab:multiseed}
\begin{tabular}{lcccc|cc}
\toprule
\textbf{Dataset} & S\,0 & S\,1 & S\,2 & S\,42 & Mean & Std \\
\midrule
IHC       & 0.806 & 0.799 & 0.797 & 0.800 & 0.800 & 0.003 \\
ISHate    & 0.896 & 0.895 & 0.903 & 0.906 & 0.900 & 0.005 \\
Vicomtech & 0.804 & 0.847 & 0.832 & 0.839 & 0.831 & 0.016 \\
\bottomrule
\end{tabular}
\end{center}

\section{Comparison with Prior Work}

\begin{center}
\captionof{table}{IHC binary comparison with ElSherief et al.~\cite{elsherief2021latent}. Evaluated on \texttt{not\_hate} vs.\ \texttt{implicit\_hate} only, matching their Table~3 setup.}
\label{tab:ihc_comparison}
\begin{tabular}{lccc}
\toprule
\textbf{Model} & \textbf{F1} & \textbf{Prec.} & \textbf{Rec.} \\
\midrule
BERT (No-index)     & 0.778 & 0.781 & 0.776 \\
HateBERT (No-index) & 0.780 & 0.785 & 0.777 \\
RoBERTa (No-index)  & 0.789 & 0.793 & 0.786 \\
\midrule
BERT (RAC)          & 0.767 & 0.767 & 0.767 \\
RoBERTa (RAC)       & \textbf{0.795} & \textbf{0.799} & \textbf{0.791} \\
\midrule
SVM n-grams~\cite{elsherief2021latent} & 0.644 & 0.614 & 0.677 \\
BERT~\cite{elsherief2021latent}        & 0.689 & 0.721 & 0.660 \\
BERT+Aug~\cite{elsherief2021latent}    & 0.704 & 0.678 & 0.732 \\
\bottomrule
\end{tabular}
\end{center}

\vspace{0.5em}

\begin{center}
\captionof{table}{ISHate Task A comparison with Ocampo et al.~\cite{ocampo2023indepth}: 3-class (Non-HS / Explicit HS / Implicit HS). RAC variants use E+K. RAC does not improve Implicit HS F1 as the binary retrieval signal conflates all three classes.}
\label{tab:ishate_taskA}
\resizebox{\textwidth}{!}{%
\begin{tabular}{lccccccccc}
\toprule
 & \multicolumn{3}{c}{\textbf{Non-HS}} & \multicolumn{3}{c}{\textbf{Explicit HS}} & \multicolumn{3}{c}{\textbf{Implicit HS}} \\
\cmidrule(lr){2-4} \cmidrule(lr){5-7} \cmidrule(lr){8-10}
\textbf{Model} & P & R & F1 & P & R & F1 & P & R & F1 \\
\midrule
BERT (No-index)     & 0.907 & 0.898 & 0.902 & 0.820 & 0.836 & 0.828 & 0.436 & 0.425 & 0.431 \\
BERT (RAC E+K)      & 0.914 & 0.623 & 0.741 & 0.588 & 0.883 & 0.706 & 0.262 & 0.403 & 0.318 \\
HateBERT (No-index) & 0.908 & 0.891 & 0.900 & 0.814 & 0.841 & 0.827 & 0.435 & 0.430 & 0.432 \\
HateBERT (RAC E+K)  & 0.918 & 0.644 & 0.757 & 0.600 & 0.879 & 0.713 & 0.253 & 0.392 & 0.308 \\
RoBERTa (No-index)  & 0.920 & 0.905 & 0.912 & 0.821 & 0.857 & 0.838 & 0.454 & 0.398 & 0.424 \\
RoBERTa (RAC E+K)   & 0.895 & 0.903 & 0.899 & 0.787 & 0.815 & 0.801 & 0.486 & 0.285 & 0.359 \\
\midrule
DeBERTa~\cite{ocampo2023indepth}      & 0.927 & 0.899 & 0.825 & 0.825 & 0.880 & 0.851 & 0.467 & 0.419 & 0.442 \\
HateBERT+ALL~\cite{ocampo2023indepth} & 0.903 & 0.896 & 0.899 & 0.827 & 0.827 & 0.827 & \textbf{0.502} & \textbf{0.559} & \textbf{0.529} \\
DeBERTa+RI~\cite{ocampo2023indepth}   & 0.922 & 0.894 & 0.908 & 0.821 & 0.878 & 0.849 & 0.460 & 0.398 & 0.427 \\
\bottomrule
\end{tabular}}
\end{center}

\section{Ablation: No-IHC Index}

\begin{center}
\captionof{table}{Ablation results on IHC when all IHC training examples are excluded from the retrieval index. The no-IHC Example index is built from \texttt{chunks\_example\_noihc.csv}. Differences relative to the standard Example index remain within seed variance ($0.800\pm0.003$).}
\label{tab:noihc_ablation}
\begin{tabular}{lccc}
\toprule
\textbf{Model} & \textbf{F1} & \textbf{Prec.} & \textbf{Rec.} \\
\midrule
BERT (E, no-IHC)     & 0.772 & 0.784 & 0.765 \\
HateBERT (E, no-IHC) & 0.787 & 0.793 & 0.783 \\
RoBERTa (E, no-IHC)  & \textbf{0.801} & \textbf{0.811} & \textbf{0.794} \\
\bottomrule
\end{tabular}
\end{center}

\end{document}

